% 来源：OEBPS/Text/01_25.xhtml（OPF spine 顺序；仅第一册）
\renewcommand{\bookchapter}{书单3}
\cleardoublepage
\BookTOCChapter[第三辑\q 经济]{第三辑\\ 经济}

\textbf{芒格说}
\par

\begin{MungerSaysBox}
{\kaishu “你必须知道重要学科的重要理论，并经常使用它们——要全部都用上，而不是只用几种。大多数人都只使用学过的一个学科的思维模型，比如说经济学，试图用一种方法来解决所有问题。你知道谚语是怎么说的：‘在手里拿着铁锤的人看来，世界就像一颗钉子。’这是处理问题的一种笨办法。”}
\par
\end{MungerSaysBox}

\cleardoublepage

% 本辑书目
% 来源：OEBPS/Text/01_26.xhtml（OPF spine 顺序；仅第一册）
\cleardoublepage
\BookEntry[戴维·S.兰德斯《国富国穷》]{戴维·S.兰德斯《国富国穷》}

\BookCover{images/书单12国富国穷.jpg}{[美] 戴维·S.兰德斯著| 门洪华译| 新华出版社2010年版}

\begin{BookSummaryBox}
\textbf{这本书回答的正是：为什么一些国家富裕，而另一些国家贫穷？}
\par
\end{BookSummaryBox}

\subsubsection*{这本书谈什么？}

这本书回答了今天人们最感兴趣、同时也争议最大的一个问题：为什么一些国家成功实现了经济发展，而另一些国家却深陷贫穷的泥淖？
\par

作者批判、发展了斯密《国富论》的基本观点，将地理、历史、科技、哲学、宗教、国际政治、文化等因素综合在一起，多层次地分析了财富创造、国家兴衰更替的根本原因，提出了人类应该吸取历史教训，创造性地走出一条人、自然与社会协调发展的新路。
\par

这本书因其主旨宏大，思想精深，而被誉为“新《国富论》”。
\par

\subsubsection*{芒格为什么谈起这本书？}

在伯克希尔·哈撒韦1999财年的年会上，芒格说：
\par

\begin{quotation}
这本书重要的书（《国富国穷》）考察的是，随着时间的推移，财富与文化会如何相互作用，两者又会对国家产生怎样的影响。一些文化区虽然具有一定优势，却从来没有成为世界强权，这本书对这一现象也给出了引人深思的理性解释。这本书还检视了不同国家兴衰各异的原因。这些原因是每一个人都应该了解的。这样我们就可以从这些挑战中学习，而不再重复那些悲惨的命运。
\par
\end{quotation}

\subsubsection*{《国富国穷》作者戴维·S.兰德斯如是说}

“相对于善心而言，我更钟爱事实。”
\par

“荷兰人的成功，集中反映了他们在工作和经商过程中所具有的、在龟兔赛跑的寓言故事中乌龟所体现的精神。丰厚的战利品和奖赏当然好，但以长远目光来看（千万不要忽视长期效益），最重要的还是从那些不起眼的、风险较低的小宗生意中获取的利润，它们会积少成多，而且从来都不会令人失望。”
\par

% 来源：OEBPS/Text/01_27.xhtml（OPF spine 顺序；仅第一册）
\cleardoublepage
\BookEntry[伯纳德·曼德维尔《蜜蜂的寓言》]{伯纳德·曼德维尔《蜜蜂的寓言》}

\BookCover{images/书单13蜜蜂的寓言.jpg}{[荷] 伯纳德·曼德维尔著| 肖聿译| 中国社会科学出版社2002年版}

\begin{BookSummaryBox}
\textbf{“用英国语言写出的最邪恶、最聪明的书”}
\par
\end{BookSummaryBox}

\subsubsection*{这本书谈什么？}

《蜜蜂的寓言》一书的核心，是道出了西方思想史上著名的曼德维尔悖论，即“私人恶德即公共利益”。从“人是自私自利的”这一观点出发，曼德维尔认为，在人品美德和商业社会之间存在着根本张力。因此，在他看来，所谓的试图以道德说教为手段，以“公共精神”为基础，建立一个美好繁荣的社会是不现实的。因为“不是人之美德和良善，而是制度性规则使道德品行不好的人也能为社会的福祉服务；而制约人们自私的恶行的绝大多数制度和规则，都是在充满罪恶和欺骗的商业行为的人类实践中产生”。
\par

曼德维尔的思想引起许多争议，该书18世纪在英国出版后，曾引起轩然大波，招致欧洲思想界乃至社会各界的批判。但尽管如此，这本书在西方思想史上的影响却还是不可小觑，影响了休谟、亚当·斯密和哈耶克等大师。
\par

\subsubsection*{芒格为什么谈起这本书？}

在芒格“论学院派经济学”的演讲中，他谈到经济学中总有一些无法避免的悖论。比如，恶行效应引发的一个问题与妒忌有关。他说：
\par

\begin{quotation}
妒忌在摩西的律法中是饱受指责的。你们可能记得希伯来人刻在墓碑上的文字：你们不能觊觎邻居的驴子，你们不能垂涎邻居的女仆，你们不能贪图……这些古代犹太人知道妒忌的人们是什么德性，也知道这些人会惹出什么麻烦。他们对妒忌真的绝不姑息，他们这么做是正确的。但曼德维尔令人信服地——反正我是信服的——证明妒忌会极大地推动消费意欲。所以呢，妒忌既是摩西十诫严厉禁止的糟糕恶行，却又是促使经济增长的驱动力。
\par
\end{quotation}

芒格提到的曼德维尔，是哲学家和讽刺作家，他在1705年发表了一首讽刺时局的诗歌：《蜜蜂的寓言，或，私人的恶行，公共的利益》。曼德维尔的哲学观是，利他主义损害了国家利益和学术进步，人类的自私心反倒是真正推动进步的引擎。因而他得到了“私心即是公益”的悖论。
\par

\subsubsection*{《蜜蜂的寓言》作者伯纳德·曼德维尔如是说}

“私人恶德即公众利益。”
\par

“若要界定作为奖赏的、表现得最充分的光荣，要廓清对光荣最多见的说法，便应当说：光荣乃是一种至高无上的快乐。一个人意识到自己做出了高尚的行为，便会在自爱当中享受这种乐，而他亦同时期待来自他人的喝彩。”
\par

“有理智、有教养者以最为娴熟的技巧掩盖自傲之时，亦是他最能取悦于自傲之心的时刻。”
\par

% 来源：OEBPS/Text/01_28.xhtml（OPF spine 顺序；仅第一册）
\cleardoublepage
\BookEntry[本杰明·M.弗里德曼《经济增长的道德意义》]{本杰明·M.弗里德曼《经济增长的道德意义》}

\BookCover{images/书单14经济增长的道德意义.jpg}{[美] 本杰明·M.弗里德曼著| 李天有译| 中国人民大学出版社2013年版}

\begin{BookSummaryBox}
\textbf{具有跨学科思想的经济理论著作}
\par
\end{BookSummaryBox}

\subsubsection*{这本书谈什么？}

对任何一个国家，经济增长都是一个古老的话题。在《经济增长的道德意义》中，弗里德曼开门见山地提出：“我们那么明确而极大地关注经济增长是否正确呢？”在做了大量细致的跨学科研究之后，弗里德曼最终指出：增长，而不仅仅是生活水平，是在第三世界政治与社会自由化的关键。但是，他也警告说，即便是像美国这样富裕的国家，一旦收入增长长期停滞，民主价值也会处于危险之中。
\par

\subsubsection*{芒格为什么谈起这本书？}

芒格在发表了“论学院派经济学”的演讲之后，看到一本书，是阿尔弗雷德·诺普夫出版社在2005年出版的。芒格说：
\par

\begin{quotation}
它的作者是杰出的哈佛大学经济学教授本杰明·弗里德曼。如我在演讲中所希望的那样，这本书讨论了经济学和道德之间的相互关系。这本书的名字是《经济增长的道德意义》。读者从书名可以看出，弗里德曼教授尤其感兴趣的是经济增长对道德的影响，而我感兴趣的则与此相反，主要是道德对经济增长的影响。这个区别关系不大，因为任何受过教育的人都能明白这两者会相互产生或好或坏的影响，造成通常所说的“良性循环”或者“恶性循环”。弗里德曼教授给这个主题增添了一句来自艾利沙·本·阿萨里亚拉比（公元前1世纪左右第二代犹太密西拿学者）的名言：“没有面包就没有法律，没有法律就没有面包。”
\par
\end{quotation}

\subsubsection*{《经济增长的道德意义》作者本杰明·M.弗里德曼如是说}

“关于经济增长的问题并不是一个物质或道德价值的问题。确实，经济增长常常伴随着不必要的副作用，如对传统文化的破坏和对环境的损害，但是这些问题中有一些确实是我们有权利加以考察的一种正当的道德关怀。虽然如此，经济增长也包含着道德上有益的社会和政治结果。特别是在评价不同的公共政策时，重要的是不仅要考虑我们熟悉的道德上的负面因素，同时也要考虑到这些道德上的正面因素。”
\par

“正确理解有关经济增长的公共政策和私人行动之间的关系也一样至关重要。这里，同样的结果是，上升的生活水准的正面道德结果会显著地改编故事的内容。”
\par

% 来源：OEBPS/Text/01_29.xhtml（OPF spine 顺序；仅第一册）
\cleardoublepage
\BookEntry[大卫·李嘉图《政治经济学及赋税原理》]{大卫·李嘉图《政治经济学及赋税原理》}

\BookCover{images/书单15政治经济学及赋税原理.jpg}{[英] 大卫·李嘉图著| 郭大力、王亚南译| 商务印书馆2013年版}

\begin{BookSummaryBox}
\textbf{英国古典政治经济学的巅峰之作}
\par
\end{BookSummaryBox}

\subsubsection*{这本书谈什么？}

《政治经济学及赋税原理》批判继承了亚当·斯密的劳动价值论，并以此为基础阐述了工资、利润和地租的相互关系，以及货币理论、对外贸易中的比较成本学说、赋税的一般原理和原则。
\par

在李嘉图看来，人类行为的准则是个人利益，是追求个人功利，即寻求快乐和避免痛苦，而追求个人利益同整个社会的利益是一致的，符合社会最大多数人的最大幸福。他实际上是站在工业资产阶级立场上，把资本家追求利润的个人利益与生产发展、社会财富增加和社会进步看作是一致的。不过，这种研究方法带有形而上学的片面性，把资本主义生产关系和制度看作是永恒的、不变的。
\par

\subsubsection*{芒格为什么谈起这本书？}

2003年10月3日，芒格在“论学院派经济学”的演讲中谈道：
\par

\begin{quotation}
如果美国和中国进行贸易，那么美国人的生活水平会得到改善，李嘉图证明了这一点。但是哪个国家的经济会发展得更快呢？那显然是中国。在自由贸易的推动之下，他们吸收了世界上各种现代科技，而且正如亚洲四小龙已经证明的，他们很快就会走到前面去。看看中国的香港和台湾，看看早年的日本就知道了。
\par
\end{quotation}

比较优势通过自由贸易带来的好处通常遭到忽略，在其1817年出版的名著《政治经济学及赋税原理》中，大卫·李嘉图对此进行了揭示：
\par

\begin{quotation}
葡萄牙生产红酒和毛呢的成本可能比英国低。然而，这两个国家生产这两种产品的相对成本是有差别的。在英国，生产红酒很难，但生产毛呢的难度并不是很大。在葡萄牙，这两者都很容易生产。因此，尽管葡萄牙生产毛呢的成本比英国低，但如果葡萄牙生产更多的红酒，用来交换英国的毛呢，那么葡萄牙的毛呢成本会更低。相反地，英国也从这种贸易中获益，因为它生产毛呢的成本没有变化，但它现在能够用更接近于毛呢的成本来获得红酒。
\par
\end{quotation}

芒格同时还指出：
\par

\begin{quotation}
国家之间相互“委派”任务能够获得比较优势，同样地，管理人员在安排工作时也可以利用李嘉图的这个原理，而这是很多人都没想到的。即使一个管理人员能够亲自更好地完成各种工作，把这些工作分散到各人手里仍然是对大家都有利的做法。
\par
\end{quotation}

\subsubsection*{《政治经济学及赋税原理》作者大卫·李嘉图如是说}

“金属和贵金属同其他物品一样，是由劳动生产的，因而其价值也是由劳动决定的，并且是随着劳动量的变动而变动的。”
\par

“货币的功能主要是流通的手段。”
\par

“当流通中的货币数量多于正常需要水平时，商品价格就上涨；相反，当流通中货币数量小于正常需要水平时，商品价格下跌。假定一国的商品流通量和价值是已知数，则‘一国所能投入的货币量必须取决于货币的价值’。”
\par

% 来源：OEBPS/Text/01_30.xhtml（OPF spine 顺序；仅第一册）
\cleardoublepage
\BookEntry[格里高利·曼昆《经济学原理》（第7版）]{格里高利·曼昆《经济学原理》（第7版）}

\BookCover{images/书单16经济学原理.jpg}{[美] 格里高利·曼昆著| 梁小民、梁砾译| 北京大学出版社2015年版}

\begin{BookSummaryBox}
\textbf{世界上最受欢迎的经济学教材}
\par
\end{BookSummaryBox}

\subsubsection*{这本书谈什么？}

哈佛大学经济学教授曼昆的《经济学原理》(第7版)是目前国内市场上最受欢迎的经济学教材，分“微观经济学”和“宏观经济学”两册。它的最大特点是“学生导向”，与同类教科书相比，本书更强调经济学原理的应用和政策分析，而不是正式的经济学模型。
\par

作者在大部分章节里都提供了案例，以说明经济学原理如何应用于现实经济学问题的分析。为了使读者懂得什么是生活中的经济学，作者在书中还提供了大量的“新闻摘录”。该书已被翻译成20多种语言。
\par

\subsubsection*{芒格为什么谈起这本书？}

2003年10月3日，芒格在“论学院派经济学”的演讲中，谈到学院派经济学的一些明显优点。例如，它生逢其时、生逢其地，自然得到了发展。所有计划经济都崩溃以后，那些自由市场经济或者半自由市场经济都蓬勃发展，这增加了经济学的名声。而关于经济学如何通过汲取其他学科的养分，进行跨学科研究，芒格说道：
\par

\begin{quotation}
在格里高利·曼昆撰写的教材中，我们可以发现，这位经济学家从其他学科吸取所需养分的本领已经非常高明了。我肯定是美国少数在那本书刚出版时就买下来的商人之一，因为那本教材得了一大笔预付稿费。我想弄清楚那个家伙到底做了些什么，怎么能够取得这么大的一笔稿费。所以我就这么凑巧把曼昆这本为大一学生写的教材给看完了。书中列举了许多经济学原理：机会成本是一种超级力量，所有希望获得正确答案的人都可以使用。还有，激励机制也是一种超级力量。
\par

最后还有“公地悲剧”的原理，这个原理是由我的老朋友，加州大学圣塔巴巴拉分校教授加勒特·哈丁提出的。哈丁为经济学引入了一只邪恶的无形之脚，它足以和斯密那只做好事的无形之手相提并论。我认为哈丁的理论使经济学变得更加完善。哈丁当年向我介绍他的理论时，我就知道他这个“公地悲剧”理论迟早会被写进教科书。你们看啊，二十年过去了，它终于被写进了经济学教材。曼昆这种借鉴其他学科、吸取哈丁的理论和其他有用知识的做法是很正确的。
\par
\end{quotation}

\subsubsection*{《经济学原理》作者格里高利·曼昆如是说}

“一个家庭面临这许多决策，它必须考虑到每个成员的能力、努力和愿望，以在其各个成员中配置稀缺资源。”
\par

“由于资源是稀缺的，社会资源的管理就显得尤为重要。”
\par

“社会上的每个人不能达到他希望的最高生活水平。”
\par

“某种东西的成本是为了得到它所放弃的东西。”
\par

“贸易可以使每个人的状况都变得更好。”
\par

“征服政策者绝不能忘记激励，因为许多政策改变了人们面临的成本和收益。”
\par

% 来源：OEBPS/Text/01_31.xhtml（OPF spine 顺序；仅第一册）
\cleardoublepage
\BookEntry[约翰·肯尼思·加尔布雷思《富裕社会》]{约翰·肯尼思·加尔布雷思《富裕社会》}

\BookCover{images/书单17富裕社会.jpg}{[加] 约翰·肯尼思·加尔布雷思著| 赵勇译| 江苏人民出版社2009年版}

\begin{BookSummaryBox}
\textbf{纽约图书馆的世纪丛书之一}
\par
\end{BookSummaryBox}

\subsubsection*{这本书谈什么？}

《富裕社会》由美国自由派经济学家加尔布雷思于50年前写就第一版，后来几经增删和再版。这本书的前几版通篇都在明确表述对通货膨胀的担心。通货膨胀一直对社会基本福利构成威胁。
\par

作者加尔布雷思笔下的“富裕社会”，可以理解为一个经济不断发展的社会，这个社会已经具有消灭贫困和流离失所的力量，只不过为了避免打击那些创造了大量财富的阶层的发展积极性，这样一个“富裕社会”漠视调节社会收入分配的重要性。
\par

在书中，加尔布雷思以“富足社会”与“传统智慧”扩展了经济学的语言，指出了现代富足社会背后隐藏的严重问题，如公共生活与私人生活标准间的资金与行动分配不均，以及环境问题的重要性。
\par

\subsubsection*{芒格为什么谈起这本书？}

2000年11月10日，芒格在慈善圆桌会议早餐会上的讲话，主要内容为美国股市上涨带来的所谓“财富效应”。美国的股票价格对经济产生了多大的影响呢？芒格认为，传统经济学家考虑这个问题所用的数据与现实有很大脱节。除了没有考虑退休金的因素外，也没有把“黑金”（bezzle）的因素考虑在内。
\par

\begin{quotation}
“黑金”这个词跟贪污有关，哈佛大学经济学系教授约翰·肯尼思·加尔布雷思用它指在尚未败露的贪污中得到的金钱。加尔布雷思发现，黑金对消费有非常强烈的刺激作用。毕竟，贪污者花钱更大手大脚，因为他的钱来得更容易，而且他的雇主的支出将会一如既往，因为雇主尚不知道其财产已经被“偷”走了。
\par
\end{quotation}

在芒格看来，加尔布雷思并没有铺开他的洞见，他满足于（在经济学思考上）提供一些牛虻式的刺激。所以芒格进一步发挥了加尔布雷思的“黑金”概念，提出“灰金”“捞灰金”和“灰金客”的概念：
\par

\begin{quotation}
正如凯恩斯指出的，在依靠劳动换取收入的原始经济中，当女裁缝把一件衣服以20美元的价格卖给鞋匠时，鞋匠就少了20美元可以消费，而女裁缝则多了20美元可以用。总消费支出并没有受到lollapalooza效应\footnote{芒格认为，几乎每个系统都受到多种因素的影响，所以要理解这样的系统，就必须熟练运用来自不同学科的多元思维模型。当几个模型联合起来，共同作用一个方向，就得到了芒格所谓的“lollapalooza效应”。——编者注}的影响。但如果政府印刷了另外一张20美元的钞票，用它来买一双鞋，鞋匠多得到了20美元，可是没有人觉得自己的钱变少了。当鞋匠下次再买一件衣服的时候，这个过程就重演了，不会无休止地持续放大，但会产生所谓的凯恩斯乘数效应，这是一种促进消费的lollapalooza效应。同样地，和同等规模的诚实交易相比，尚未败露的贪污得来的钱对消费的刺激效应更大。加尔布雷思是苏格兰人，喜欢深刻地揭示生活的世态炎凉。毕竟，这个苏格兰人还热衷于接受命中注定、无法改变的婴儿诅咒这样的荒唐想法。我们大多数人并不喜欢加尔布雷思的观点。但我们不得不承认，他有关“黑金”的看法基本上是正确的。
\par
\end{quotation}

加尔布雷思出生在加拿大安大略省，从安大略农学院毕业后，到伯克利加利福尼亚大学攻读博士学位。他于1949年到哈佛大学经济学系担任教职，是肯尼迪总统的好友，1961年到1963年担任美国驻印度大使。作为一名经济学家，加尔布雷思持有改良派的价值观，他写了许多描述经济理论如何与现实生活相悖的著作。《富裕社会》便是他的代表作之一。
\par

\subsubsection*{《富裕社会》作者约翰·肯尼思·加尔布雷思如是说}

“我们都认为悲观是智力超群的标志。”
\par

“美当然没有绝对的标准。所以追求美才会如此有趣。”
\par

“如果你什么都不想做，那么你就去开会吧。”
\par

“经济预测的唯一功能是让占星学变得令人尊敬。”
\par

“现代的保守主义者从事的是道德哲学中最古老的活动，也就是说，他们要为自私寻找一个超级漂亮的道德借口。”
\par

“许多美国人死于食物过多，而非死于食物太少。”
\par

“当面临要么改变想法，要么证明无须这么做的选择时，绝大多数人都会忙于寻找证据。那些走运的人总是说他们过得幸福是因为他们的道德很高尚。”
\par

“难道你不明白吗？让一家企业停泊在废话当中，和让它驶入麻烦的思维之海比起来，是一种好得太多、太多的做法。”
\par

“有权势的人总是宁愿拿全副身家去冒险，也不肯舍弃一点利益。”
\par

“在任何大型组织里面，和多数人一起错比一个人对安全得多。”
\par

% 来源：OEBPS/Text/01_32.xhtml（OPF spine 顺序；仅第一册）
\cleardoublepage
\BookEntry[亚当·斯密《国民财富的性质和原因的研究》]{亚当·斯密《国民财富的性质和原因的研究》}

\BookCover{images/书单18国民财富的性质和原因的研究.jpg}{[英] 亚当·斯密著| 郭大力、王亚南译| 商务印书馆1974年版}

\begin{BookSummaryBox}
{\centering \textbf{第一部系统的经济学经典\\ 影响人类文化进程的100部经典之一}\par}
\end{BookSummaryBox}

\subsubsection*{这本书谈什么？}

1776年，《国民财富的性质和原因的研究》（简称《国富论》）出版，这是经济学之父亚当·斯密在亲赴欧洲各地考察之后，撰写的一部改变人类历史的创世巨著。
\par

那正是英国资本主义的成长时期，英国手工制造业正在开始向大工业过渡，而英国产业的发展在很大的程度上还受着封建制度残余和流行一时的重商主义政策的束缚。年轻的英国资产阶级，为了扫除前进道路上的障碍，正迫切要求一个自由的经济学说体系，为它鸣锣开道。亚当·斯密的《国富论》，就是在这个历史时期，肩负着有这样的阶级历史任务而问世了。这部书总结了近代初期各国资本主义发展的经验，并在批判吸收了当时有关重要经济理论的基础上，就整个国民经济运动过程作了较系统、较明白的描述。它不但对于英国资本主义的发展，直接产生了重大的促进作用，而且对世界资本主义的发展来说，恐怕也没有任何其他一部资产阶级的经济学著作，曾产生那么广泛的影响。
\par

\subsubsection*{芒格为什么谈起这本书？}

芒格在斯坦福大学法学院的演讲中，提出了自己对真正的学术观念的看法，他说道：
\par

\begin{quotation}
所以我认为，学术蕴含了许多普世智慧，而且最好的学术观念确实是有用的……当然，当我谈到跨学科方法——你们应该掌握各个学科的主要模型，并将它们派上用场——时，我是真的呼吁你们不要理会学科的法定界限……你们不需要了解所有的知识，只要吸取各个学科最杰出的思想就行了。
\par
\end{quotation}

芒格还常常提到“与已逝的伟人”交朋友，把重要的概念与重要的任务联系在一起，会帮助你理解更多。当然，他还特别说道：
\par

\begin{quotation}
如果你能与亚当·斯密交朋友，那你的经济学一定学得非常好。
\par
\end{quotation}

\subsubsection*{《国民财富的性质和原因的研究》作者亚当·斯密如是说}

“我们会经常发现，在世界上具有伟大人性但是却缺乏自我控制的人，在追求最高荣誉时，一旦碰到困难和危险，就懒惰，犹豫，容易沮丧；相反，我们也常常发现能够完善地进行自我控制的人，没有任何困难可以吓倒他们，没有任何危险能够惊骇他们，但同时，他们对有关正义或人性的全部感觉似乎无动于衷。”
\par

“科学是解狂热和迷信之毒的良药。”
\par

“劳动技能的影响要大于有用劳动占比的影响，文明社会里产品的数量更多可以表明这一点。”
\par

“我们不能借着肉贩、啤酒商或面包师的善行而获得晚餐，而是源于他们对自我利益的看重。”
\par

